%% ============================================================
%%  chapitre3.tex — الفصل الثالث: الحالات التي يسمح فيها القانون
%%                  باللجوء إلى الشغل المؤقت
%% ============================================================

\chapter{الحالات التي يسمح فيها القانون باللجوء إلى الشغل المؤقت}
\label{chap:chapitre3}

%% ─── مقدمة الفصل ──────────────────────────────────────────
\section*{مقدمة الفصل}
\phantomsection\addcontentsline{toc}{section}{مقدمة الفصل}

الأصل في مدونة الشغل أن تقوم علاقة الشغل على الاستقرار، وأن يكون العقد غير محدد المدة (\textenglish{CDI}) هو القاعدة. ولذلك لم يُجز المشرع اللجوء إلى الأجير المؤقت إلا استثناءً، وفي أربع حالات أوردتها المادة \textenglish{496} على سبيل الحصر، لا يجوز التوسع فيها ولا القياس عليها. ويُشترط فضلاً عن ذلك أن تستشير المقاولة المستعملة الهيئات التمثيلية لأجرائها قبل اللجوء إلى هذه الصيغة.

ويجيب هذا الفصل عن سؤالين محوريين من أسئلة النقاش في النازلة: \textbf{السؤال الثالث} (ما الحالات التي يسمح فيها القانون باللجوء إلى شركات التشغيل المؤقت؟) و\textbf{السؤال الخامس} (هل يمكن استعمال الشغل المؤقت لتلبية حاجيات دائمة ومستمرة للمقاولة؟)، عبر أربعة محاور متوازنة:
\begin{itemize}[leftmargin=*, itemsep=4pt]
    \item \textbf{أولاً:} استبدال أجير توقف عقده؛
    \item \textbf{ثانياً:} ارتفاع مؤقت في نشاط المقاولة؛
    \item \textbf{ثالثاً:} الأشغال الموسمية والأشغال التي استقر العرف بشأنها (الظرفية)؛
    \item \textbf{رابعاً:} الحالات الممنوعة ومدى جواز تلبية الحاجيات الدائمة وإسقاطها على المجموعة \textenglish{8}.
\end{itemize}

\begin{table}[htbp]
\centering
\small
\renewcommand{\arraystretch}{1.5}
\begin{tabular}{|p{4.4cm}|p{4.6cm}|p{4.4cm}|}
\hline
\textbf{حالة اللجوء (المادة \textenglish{496})} & \textbf{الشرط الأساسي} & \textbf{المدة القصوى (المادة \textenglish{500})} \\
\hline
استبدال أجير توقف عقده & توقف مؤقت لعقد أجير محدد، لغير سبب الإضراب & مدة توقف عقد الأجير المعوَّض \\
\hline
ارتفاع مؤقت في النشاط & زيادة فعلية وعابرة في حجم العمل & \textenglish{3} أشهر قابلة للتجديد مرة واحدة \\
\hline
الأشغال الموسمية & نشاط يتكرر في فترات محددة من السنة & \textenglish{6} أشهر غير قابلة للتجديد \\
\hline
الأشغال التي جرى العرف بعدم إبرام عقود غير محددة المدة بشأنها & طبيعة القطاع أو النشاط تجعل العمل مؤقتاً بطبيعته & \textenglish{6} أشهر غير قابلة للتجديد \\
\hline
\end{tabular}
\caption{الحالات الحصرية للجوء إلى الشغل المؤقت وشروطها ومددها}
\label{tab:cas-recours}
\end{table}

%% ─────────────────────────────────────────────────────────
\section{استبدال أجير توقف عقده}

\subsection{مضمون الحالة ونطاق التوقف المشروع}
تسمح المادة \textenglish{496} بوضع أجير مؤقت رهن إشارة المقاولة المستعملة \textbf{لتعويض أجير غائب أو توقف عقد شغله}. ويقتضي ذلك أن يكون عقد الأجير الأصلي قائماً لكنه معلَّق مؤقتاً، بحيث يعود صاحبه إلى منصبه عند زوال سبب التوقف. ومن أبرز صور هذا التوقف المشروع:
\begin{itemize}[leftmargin=*, itemsep=4pt]
    \item الغياب بسبب المرض غير المهني أو حادثة الشغل أو المرض المهني.
    \item إجازة الولادة أو التمديد المرتبط بها.
    \item العطلة السنوية المؤدى عنها أو التغيبات المرخص بها قانوناً أو اتفاقياً.
    \item الإيقاف المؤقت للعقد لأي سبب آخر يقره القانون (كالخدمة أو التفرغ النقابي المرخص).
\end{itemize}

\subsection{شروطها وحدودها القانونية}
\begin{enumerate}[leftmargin=*, itemsep=6pt]
    \item \textbf{استثناء الإضراب:} لا يجوز اللجوء إلى الأجير المؤقت إذا كان توقف العقد ناتجاً عن ممارسة حق الإضراب (المادة \textenglish{496})، حتى لا يُستعمل الشغل المؤقت أداةً لكسر الإضرابات وإفراغ هذا الحق الدستوري من محتواه.
    \item \textbf{ارتباط المدة بالغياب:} لا تتجاوز مدة المهمة مدة توقف عقد الأجير المعوَّض (المادة \textenglish{500})، فإذا عاد هذا الأخير انتهى مبرر المهمة فوراً.
    \item \textbf{تحديد الأجير المعوَّض:} يجب أن يُذكر سبب اللجوء في العقد المكتوب بين مقاولة التشغيل المؤقت والمقاولة المستعملة (المادة \textenglish{499})، مما يقتضي عملياً تحديد الأجير الذي يُعوَّض وسبب غيابه، حتى يتسنى لمفتش الشغل التحقق من صحة السبب.
\end{enumerate}

%% ─────────────────────────────────────────────────────────
\section{ارتفاع مؤقت في النشاط}

\subsection{مضمون التزايد المؤقت ومعايير تحققه}
تسمح المادة \textenglish{496} باللجوء إلى الأجير المؤقت لمواجهة \textbf{تزايد مؤقت في نشاط المقاولة}، كطلبية استثنائية، أو ذروة غير معتادة في الإنتاج، أو مشروع عابر يتجاوز طاقة الأجراء الدائمين. والمعيار هنا مزدوج: أن تكون الزيادة \textbf{حقيقية} يمكن إثباتها محاسبياً وتجارياً، وأن تكون \textbf{مؤقتة} بطبيعتها، أي أن ينتظر زوالها في أجل منظور، لا أن تمثل توسعاً دائماً في حجم الإنتاج.

\subsection{القيود الزمنية والمنع بعد الفصل الاقتصادي}
\begin{enumerate}[leftmargin=*, itemsep=6pt]
    \item \textbf{مدة محدودة:} ثلاثة أشهر قابلة للتجديد \textbf{مرة واحدة فقط} (المادة \textenglish{500})، أي ستة أشهر كحد أقصى. فإذا استمرت «الزيادة» بعد ذلك، فهي في الحقيقة مستوى عادي من النشاط يستوجب تشغيلاً دائماً.
    \item \textbf{منع اللجوء بعد الفصل لأسباب اقتصادية:} لا يجوز للمقاولة التي قامت بفصل أجراء لأسباب اقتصادية أن تلجأ إلى أجراء مؤقتين لمواجهة تزايد مؤقت في النشاط، وذلك خلال السنة الموالية للفصل وبالنسبة للمناصب التي شملها (المادة \textenglish{498})، مع مراعاة أحكام الأولوية في إعادة التشغيل وفق المادة \textenglish{508}.
    \item \textbf{عدم التحايل بالتعاقب:} لا يجوز تحويل هذه الحالة إلى غطاء دائم بتعاقب مهام قصيرة على المنصب نفسه أو بتدوير الأجراء، لأن العبرة بطبيعة الحاجة لا بشكل العقود المتتالية.
\end{enumerate}

%% ─────────────────────────────────────────────────────────
\section{الأشغال الموسمية والظرفية}

تجمع هذه الفقرة حالتين أوردتهما المادة \textenglish{496} تشتركان في أن العمل فيهما مؤقت بطبيعته الموضوعية، لا بقرار إداري من المشغل.

\subsection{مضمون الأشغال الموسمية والأشغال التي استقر العرف بشأنها}
\begin{itemize}[leftmargin=*, itemsep=6pt]
    \item \textbf{الأشغال الموسمية:} هي الأشغال التي تتكرر في فترات محددة من السنة تبعاً لإيقاع طبيعي أو اقتصادي، كجني المحاصيل الفلاحية وتلفيفها، أو ذروة الموسم السياحي، أو بعض الصناعات الغذائية المرتبطة بمواسم الإنتاج. ويقوم معيارها على التكرار الدوري للنشاط وانقطاعه خارج الموسم.
    \item \textbf{الأشغال التي جرى العرف بعدم إبرام عقود غير محددة المدة بشأنها (الأشغال الظرفية):} وهي أنشطة يكون العمل فيها بالمهمة أو بالمشروع عرفاً مستقراً في القطاع بالنظر إلى طبيعتها المؤقتة، كالإنتاج السمعي البصري والسينمائي، وتنظيم التظاهرات الثقافية والرياضية والمعارض.
\end{itemize}

\subsection{النظام الزمني الصارم (ستة أشهر غير قابلة للتجديد)}
تحدد المادة \textenglish{500} مدة المهمة في هاتين الحالتين في \textbf{ستة أشهر غير قابلة للتجديد}. ويعكس هذا السقف الصارم منطق المشرع: فالموسم أو التظاهرة الظرفية محدودة بطبيعتها في الزمان، وما يتجاوز ستة أشهر لم يعد موسماً أو نشاطاً عرضياً بل عملاً مستمراً يفرض اللجوء إلى التعاقد الدائم.

%% ─────────────────────────────────────────────────────────
\section{الحالات الممنوعة (الحاجيات الدائمة والمستمرة)}

\subsection{الموانع التشريعية الصريحة ومدى جواز تلبية الحاجيات الدائمة}
أقر المشرع جملة من الموانع الآمرة التي يُحظر فيها اللجوء إلى الشغل المؤقت:
\begin{tcolorbox}[colback=red!5, colframe=accentcolor, arc=4pt,
    title={\textbf{حالات يُمنع فيها اللجوء إلى الشغل المؤقت}}]
\begin{itemize}[leftmargin=*, itemsep=4pt]
    \item \textbf{تعويض الأجراء المضربين} (المادة \textenglish{496}).
    \item \textbf{الأشغال التي تكتسي خطورة خاصة} (المادة \textenglish{497}).
    \item \textbf{مواجهة تزايد مؤقت في النشاط داخل المناصب التي شملها فصل لأسباب اقتصادية}، خلال السنة الموالية للفصل (المادة \textenglish{498}).
    \item \textbf{كل لجوء خارج الحالات الأربع للمادة \textenglish{496}}، أو دون استشارة الهيئات التمثيلية للأجراء.
\end{itemize}
\end{tcolorbox}

وبشأن التساؤل حول \textbf{إمكانية استعمال الشغل المؤقت لتلبية حاجيات دائمة ومستمرة للمقاولة}، فإن الجواب هو \textbf{بالنفي القاطع}، وإن لم تصُغه مدونة الشغل في جملة منفردة على غرار المادة \textfrench{L. 1251-5} من قانون العمل الفرنسي. ويُستخلص هذا المنع في القانون المغربي من ثلاثة أسس متكاملة:
\begin{enumerate}[leftmargin=*, itemsep=5pt]
    \item \textbf{الطابع الحصري للحالات:} كل حالات المادة \textenglish{496} مبنية على فكرة الحاجة العابرة، ولا مكان بينها للحاجة الدائمة.
    \item \textbf{المدد القصوى:} لا تتجاوز المهمة في أي حالة ستة أشهر (المادة \textenglish{500}) باستثناء الإحلال المرتبط بمدة غياب محدد، ومن ثم يستحيل قانوناً أن يشغل أجير مؤقت منصباً قاراً لسنوات.
    \item \textbf{مبدأ استقرار الشغل:} باعتبار العقد غير محدد المدة هو الأصل العام، يُفسَّر كل استثناء تفسيراً ضيقاً. ويترتب عن خرق هذا المنع جزاء زجري في المادة \textenglish{505} (غرامة \textenglish{2.000} إلى \textenglish{5.000} درهم)، وجزاء مدني يستخلصه القضاء بإعادة التكييف إلى عقد غير محدد المدة.
\end{enumerate}

\subsection{إسقاط الضوابط الحصرية على وضعية المجموعة 8}
في الوضعية المدروسة للمجموعة \textenglish{8}، يشتغل العمال الذين توفرهم شركة المناولة داخل المصنع \textbf{منذ سنوات}، ويؤدّون المهام ذاتها التي يقوم بها أجراء الشركة الأصلية. ولا تندرج هذه الوضعية في أي من الحالات الأربع للمادة \textenglish{496}:
\begin{itemize}[leftmargin=*, itemsep=4pt]
    \item ليس هناك أجير غائب يُعوَّض، بل مناصب قارة في خط الإنتاج.
    \item لا يمكن الحديث عن زيادة مؤقتة في النشاط تستمر لسنوات، والحد الأقصى هو ستة أشهر.
    \item لا يتعلق الأمر بنشاط موسمي ولا بقطاع جرى عرفه على العمل المؤقت.
\end{itemize}
فالأمر يتعلق إذن بتلبية حاجيات دائمة ومستمرة، وهو ما يحظره القانون. ويفسر ذلك لجوء بعض المقاولات إلى تسمية العقد «مناولة» بدل «شغل مؤقت»، غير أن العبرة في القضاء بحقيقة العلاقة لا بتسميتها.

%% ─── خلاصة الفصل ──────────────────────────────────────────
\section*{خلاصة الفصل}
\phantomsection\addcontentsline{toc}{section}{خلاصة الفصل}

حصر المشرع المغربي اللجوء إلى الشغل المؤقت في أربع حالات تجمعها فكرة الحاجة العابرة، وأحاطها بمدد قصوى لا تتجاوز ستة أشهر في الغالب، وبموانع صريحة تتعلق بالإضراب والأشغال الخطيرة والفصل الاقتصادي. ويُستخلص أن الشغل المؤقت لا يمكن قانوناً أن يُستعمل لتلبية حاجيات دائمة ومستمرة، وهو ما يجعل وضعية عمال المجموعة \textenglish{8} خارج كل إطار مشروع. ويبقى بعد ذلك تحديد الالتزامات الملقاة على مقاولة التشغيل المؤقت والمقاولة المستعملة، وهو موضوع الفصل الرابع.
